\chapter{OpenFlow: Enabling Innovation in Campus Networks}\label{sec:openflow}
\paperinfo{N. McKeown, T. Anderson, H. Balakrishnan, G. Parulkar, L. Peterson, J. Rexford, S. Shenker, J. Turner}{ACM SIGCOMM CCR, Vol.\ 38 No.\ 2, April 2008}

\begin{keybox}[OpenFlow in one paragraph]
\key{OpenFlow is based on an Ethernet switch, with an internal flow table and a standardized interface to add and remove flow entries.} The goal is for vendors to add OpenFlow to switches deployed in college campus backbones and wiring closets, so researchers can \key{run experiments} on real networks at line rate, while \key{vendors do not need to expose the internal workings of their switches}. It is a pragmatic compromise between full programmability and what vendors will accept.
\end{keybox}

\section{The need for programmable networks}
\begin{itemize}
\item The huge installed base of equipment and protocols, and reluctance to experiment on production traffic, mean that \key{today there is almost no practical way to experiment with new network protocols} (new routing protocols, alternatives to IP) in realistic settings. \key{New ideas go untried and untested}, hence the belief that the network infrastructure has \key{``ossified''}.
\item \key{GENI} is a proposed \key{nationwide research facility for experimenting with new network architectures and distributed systems}. Such programmable networks need \key{programmable switches and routers} that, using virtualization, process packets for several isolated experimental networks at once. In GENI a \key{researcher is allocated a \emph{slice} of resources across the whole network}: a portion of network links, packet-processing elements (e.g., routers) and end-hosts.
\item But the plans for nationwide facilities are \key{ambitious (and costly), and will take years to deploy}. The paper asks a shorter-term question: \emph{as researchers, how can we run experiments in our campus networks?}
\item Goal: \key{propose a new switch feature that extends programmability into the wiring closet of college campuses}.
\end{itemize}
\marginnote{Ossified = anything new is difficult to introduce.}

\section{Why existing options do not work}
\begin{itemize}
\item \key{Commercial switches/routers} (approach \emph{not} taken: persuading vendors to open up). They \key{do not provide an open software platform}, let alone a way to \key{virtualize their hardware or software}. External interfaces are narrow (just packet forwarding) and all internal flexibility is hidden. \key{Internals differ from vendor to vendor}, so there is no standard platform, and \key{vendors are nervous about opening interfaces inside their boxes}: they have tuned fragile distributed protocols for years, fear experiments will crash networks, and open platforms lower the barrier for competitors.
\item \key{PC-based software routers.} \key{The simplest example is a PC with several network interfaces and an OS}; all OSes route packets, and the OS can be modified to \key{process packets in almost any manner} (e.g., with Click). The \key{problem is performance}: \key{a PC supports neither the number of ports needed for a wiring closet (100+) nor the packet-processing performance} (wiring-closet switches process over 100\,Gb/s; a PC struggles to exceed 1\,Gb/s, and the gap is widening).
\item \key{Specialized hardware for line-rate processing} is \key{not quite suitable for college wiring closets}:
\begin{itemize}
\item The ATCA-based \key{Supercharged PlanetLab Platform} (a virtualized programmable router with network processors) is \key{promising long-term, but targets large switching centres and is too expensive} for wide deployment in wiring closets.
\item \key{NetFPGA} is a \key{low-cost PCI card with a user-programmable FPGA for processing packets and four Gigabit Ethernet ports} --- \key{limited to just four network interfaces}, insufficient for a wiring closet.
\end{itemize}
\end{itemize}
\textbf{Conclusion:} \key{commercial solutions are too closed} and inflexible; research solutions have too little performance or fanout, or are too expensive.
\marginnote{We are not taking the help of outside vendors (to open their platforms). Discussing currently available solutions. Specialized routers: not widespread-able --- costly! \textbf{FPGA} = Field-Programmable Gate Array: it can be reprogrammed by the customer after manufacturing. \textbf{PCI cards}: plug into the motherboard and provide additional functionality.}

\textbf{A compromise.} Rather than full generality, seek a degree of switch flexibility that is:
\begin{itemize}
\item amenable to \key{high-performance and low-cost implementations};
\item capable of \key{supporting a broad range of research};
\item assured to \key{isolate experimental traffic from production traffic};
\item \key{consistent with vendors' need for closed platforms}.
\end{itemize}

\section{The OpenFlow switch}
\subsection{Basic idea}
\begin{itemize}
\item \key{Exploit the fact that most modern Ethernet switches and routers contain flow tables} (typically built from TCAMs) that run at line rate to \key{implement firewalls, NAT, QoS, and to collect statistics}.
\item \key{Each vendor's flow table is different}, but there is a \key{common set of functions that run in many switches and routers}. OpenFlow exploits this common set.
\item OpenFlow \key{provides an open protocol to program the flow table in different switches and routers}. A network administrator can \key{partition traffic into production and research flows}. \key{Researchers control their own flows} --- choosing the routes their packets follow and the processing they receive --- so they can try new routing protocols, security models, addressing schemes, even alternatives to IP. \key{On the same network, the production traffic is isolated and processed the same way as today.}
\item The datapath of an OpenFlow switch is a \key{flow table, with an action associated with each flow entry}. The set of actions is extensible, but a minimum set is required of all switches. High performance at low cost means forgoing arbitrary per-packet handling in favour of a limited but useful set of actions.
\end{itemize}
\marginnote{An Ethernet switch works at Layer 2; a flow table works at specific L2/L3/L4 fields --- similar to a routing table but more flexible and more detailed, i.e., cross-layer (it works across layers). Actions can be extended!}

\subsection{Three parts of an OpenFlow switch}
\paperfig[0.48]{of_fig1_switch}{Idealized OpenFlow switch. The Flow Table is controlled by a remote controller via the Secure Channel using the OpenFlow Protocol (over SSL). The Secure Channel is software (sw); the Flow Table is hardware (hw).}
An OpenFlow switch consists of at least three parts:
\begin{enumerate}
\item A \key{Flow Table}, with \key{an action associated with each flow entry}, to tell the switch how to process the flow.
\item A \key{Secure Channel} that connects the switch to a remote control process (the \key{controller}), allowing commands and packets to be sent between them.
\item The \key{OpenFlow Protocol}, \key{an open and standard way for a controller to communicate with a switch}. Because of this standard interface, \key{entries in the Flow Table can be defined externally}, and researchers need not program the switch itself.
\end{enumerate}
\marginnote{Secure Channel = communication between switch and controller.}

\subsection{Two kinds of OpenFlow switch}
\begin{itemize}
\item \key{Dedicated OpenFlow switch}: a dumb \key{datapath element that forwards packets between ports}, as defined by a remote control process. It does not support normal Layer 2/Layer 3 processing.
\item \key{OpenFlow-enabled switch}: a commercial switch, router or access point enhanced with the Flow Table, Secure Channel and OpenFlow Protocol. The Flow Table typically \emph{reuses existing hardware} such as a TCAM, and the Secure Channel and protocol are ported to the switch's operating system.
\end{itemize}

\textbf{What is a flow?} Flows are defined broadly, limited only by the Flow Table implementation. \key{A flow could be a TCP connection, all packets from a particular MAC or IP address, all packets with the same VLAN tag, or all packets from the same switch port}. For non-IPv4 experiments, a flow could be all packets matching a specific (non-standard) header.

\section{Flow entries and actions}
\subsection{Fields of a flow entry}
\key{An entry in the Flow Table has three fields}:
\begin{enumerate}
\item a \key{packet header} that defines the flow;
\item the \key{action}, which defines how the packets should be processed;
\item \key{statistics}, which \key{track the number of packets and bytes for each flow, and the time since the last packet matched} (to help remove inactive flows).
\end{enumerate}

\subsection{The Type 0 header: a 10-tuple}
In \key{first-generation ``Type 0'' switches}, the flow header is the 10-tuple below.
\begin{center}
\begin{tabular}{|c|c|ccc|ccc|cc|}
\hline
In & VLAN & \multicolumn{3}{c|}{Ethernet} & \multicolumn{3}{c|}{IP} & \multicolumn{2}{c|}{TCP} \\
Port & ID & SA & DA & Type & SA & DA & Proto & Src & Dst \\
\hline
\end{tabular}
\end{center}
A TCP flow can be specified using all ten fields; an IP flow might leave out the transport ports. Each field can be a \key{wildcard}, allowing aggregation: e.g., a flow that specifies only the VLAN ID applies to all traffic on that VLAN.
\marginnote{10-tuple flow header. Some flows use all the entries while some don't (wildcards).}

\subsection{The three basic actions (dedicated switch)}
\begin{enumerate}
\item \key{Forward this flow's packets to a given port (or ports).} Routes packets through the network; expected at line rate.
\item \key{Encapsulate and forward this flow's packets to a controller.} The packet goes over the \key{Secure Channel}. Typically \key{used for the first packet in a new flow, so the controller can decide whether the flow should be added to the Flow Table}. Some experiments send \emph{all} packets to the controller.
\item \key{Drop this flow's packets.} \key{Used for security}, to curb \key{denial-of-service attacks}, or to reduce spurious broadcast discovery traffic from end-hosts.
\end{enumerate}
\marginnote{3 simple actions!}

\subsection{OpenFlow-enabled switches: isolating experimental traffic}
To win administrators' confidence, \key{OpenFlow-enabled switches must isolate experimental traffic} (processed by the Flow Table) from production traffic (processed by the normal L2/L3 pipeline). Two ways:
\begin{enumerate}
\item \key{Add a fourth action}: \key{forward this flow's packets through the switch's normal processing pipeline}.
\item \key{Define separate sets of VLANs for experimental and production traffic}.
\end{enumerate}
Every OpenFlow-enabled switch must support one approach; some support both. All Flow Tables in a network can be managed by the same controller, but the \key{OpenFlow Protocol allows a switch to be controlled by two or more controllers for increased performance or robustness}.
\marginnote{VLAN = the concept of grouping: although hosts are on the same physical network, VLANs divide them into separate logical networks (here: experimental vs.\ production). More than one controller is possible.}

\subsection{Additional features: Type 0 and Type 1}
\begin{itemize}
\item \key{If a switch supports the header formats and the four basic actions}, it is a \key{``Type 0'' switch}.
\item \key{Many switches will support additional actions}, e.g., to \key{rewrite portions of the packet header} (for NAT, or to obfuscate addresses on intermediate links) and to \key{map packets to a priority class}. Some Flow Tables will match arbitrary header fields, enabling non-IP experiments. As a particular set of such features emerges, it will be defined as a \key{``Type 1'' switch}.
\end{itemize}
\marginnote{Type 1 = a particular set of extra features.}

\section{Controllers}
\begin{itemize}
\item \key{A controller adds and removes flow entries from the Flow Table on behalf of experiments.}
\item A \key{static controller} might be a \key{simple application running on a PC} that statically connects a set of test machines for an experiment. The flows then resemble VLANs, isolating experimental traffic --- so \emph{OpenFlow is a generalization of VLANs}.
\item \key{Sophisticated controllers} \key{dynamically add/remove flows} as an experiment progresses. One \key{might support multiple researchers, each with different accounts and permissions}, letting them \key{run multiple independent experiments} on different sets of flows. Flows belonging to a researcher (e.g., identified by a policy table in the controller) can be delivered to that researcher's user-level program, which decides whether to add a new flow entry.
\item \textbf{Performance questions:} Can a central controller process new flows fast enough? What happens if it fails? Ethane's experience suggested a low-cost PC controller can handle over 10,000 new flows/second --- enough for a large campus. Scalability and redundancy are possible by making a controller (and experiments) \emph{stateless}, allowing load balancing across several devices.
\end{itemize}
\paperfig[0.5]{of_fig2_network}{A network of OpenFlow switches and OpenFlow-enabled commercial switches/access points, all managed by one controller. The commercial switch keeps its normal software and datapath next to the Secure Channel and Flow Table.}

\section{Using OpenFlow}
\subsection{Example: Amy-OSPF}
Amy has invented \emph{Amy-OSPF} to replace OSPF and wants to test it on OpenFlow switches without changing any end-host software. Amy-OSPF runs in a controller.
\begin{enumerate}
\item She defines a flow for all traffic entering the OpenFlow switch through the port her PC is connected to, with action ``encapsulate and forward all packets to a controller''.
\item When a new application flow starts, its packet reaches the controller; Amy-OSPF picks a route through the switches and adds a flow entry in every switch along the path.
\item Subsequent packets are processed quickly, at line rate, by the Flow Tables.
\end{enumerate}
Only traffic from her own desktop is affected, so she does not disrupt others.

\subsection{Experiments in a production network: two properties}
\begin{enumerate}
\item Packets of users other than Amy should be routed by a standard, tested routing protocol running in the switch from a ``name-brand'' vendor.
\item Amy should only be able to add flow entries for her own traffic, or traffic her administrator allows her to control.
\end{enumerate}
\begin{itemize}
\item \key{Property 1 is achieved by OpenFlow-enabled switches.} In Amy's experiment, the \key{default action for all packets that don't come from Amy's PC is to forward them through the normal processing pipeline}; her packets go directly to the outgoing port.
\item \key{Property 2 depends on the controller.} The controller is a platform for experiments, and \key{Property 2 can be achieved with appropriate permissions or other ways to limit researchers' power to control flow entries}. The exact mechanism depends on the controller (e.g., NOX, a follow-on to Ethane; or extended GENI management software).
\end{itemize}

\section{More example applications}
\paragraph{Example 1: Network management and access control (Ethane).} An OpenFlow switch is a generalization of Ethane's datapath switch. Ethane lets managers define a network-wide policy in a central controller, enforced by admission-control decisions for each new flow (e.g., ``guests may use HTTP only through a web proxy'', ``VoIP phones may not talk to laptops''). The controller binds packets to senders by managing name--address bindings --- it essentially takes over DNS and DHCP, authenticates users as they join, and tracks which switch port or access point they use.

\paragraph{Example 2: VLANs.} OpenFlow can give users isolated networks, as VLANs do. Static approach: declare flows specifying which ports are reachable by traffic with a given VLAN ID; traffic from one user (identified by switch port or MAC) is tagged by the switch, via an action, with the right VLAN ID. Dynamic approach: a controller authenticates users and uses their locations to tag traffic at runtime.
\marginnote{Use of VLAN ID: hosts in the same group get the same VLAN ID.}

\paragraph{Example 3: Mobile wireless VoIP clients.} In an experiment with a new call-handoff mechanism for WiFi phones, \key{VoIP clients establish a new connection over the OpenFlow-enabled network}. A \key{controller tracks the location of clients and re-routes connections} by reprogramming the Flow Tables as users move, \key{allowing seamless handoff from one access point to another}.
\marginnote{VoIP = Voice over Internet Protocol (calls over WiFi). Seamless handoff is possible with OpenFlow!}

\paragraph{Example 4: A non-IP network.} OpenFlow does not require IP, as long as the Flow Table can match the header. \key{There are several ways an OpenFlow-enabled switch can support non-IP traffic}: \key{flows could be identified by their Ethernet header} (MAC src/dst), by a \key{new EtherType value}, or, at the IP level, by a \key{new IP Version number}. Future switches might let a controller define a generic mask (offset + value + mask).

\paragraph{Example 5: Processing packets rather than flows.} So far the controller decides when a flow starts. Some experiments must process \emph{every} packet: an intrusion-detection system inspecting each packet, an explicit congestion-control mechanism, or converting packets between protocol formats. Two ways:
\begin{enumerate}
\item \key{Force all of a flow's packets to pass through a controller}: the \key{controller doesn't add a new flow entry into the Flow Table, so the switch defaults to forwarding every packet to the controller}. Flexible, but \key{at the cost of performance}; good for testing functionality, unlikely to suit a large deployment.
\item \key{Route them to a programmable switch that does packet processing}, e.g., a \key{NetFPGA-based programmable router}. \key{Packets can be processed at line rate in a user-definable way}. The \key{OpenFlow-enabled switch acts as a patch panel to allow the packets to reach the NetFPGA}, which may sit in the wiring closet or (more likely) in a lab.
\end{enumerate}
\paperfig[0.48]{of_fig3_netfpga}{Processing packets through an external line-rate device: the OpenFlow-enabled switch acts as a patch panel and steers selected flows (red dotted) to a NetFPGA router in a laboratory.}
\marginnote{How to process packets rather than flows?}

\section{OpenFlow Consortium and deployment}
\key{The OpenFlow Consortium aims to popularize OpenFlow and maintain the OpenFlow Switch Specification.} It is a group of university researchers and network administrators; membership is free for schools and government agencies but excludes employees of switch/router/AP vendors, to keep it free of vendor influence. A commercial switch claiming to be OpenFlow-enabled must meet the Type 0 requirements. Two Stanford buildings were being deployed with production and experimental traffic on the same OpenFlow-enabled switches, isolated by VLANs.

\section{Main takeaways}
\begin{itemize}
\item \textbf{Problem:} networks are ossified; commercial gear is closed, PCs lack ports and performance, research platforms are too costly or have too few ports; GENI is years away.
\item \textbf{Idea:} standardize access to the common flow-table functions already in most switches, without opening vendor internals.
\item \textbf{Switch:} Flow Table (hw) + Secure Channel (sw) + OpenFlow Protocol, driven by a remote controller.
\item \textbf{Flow entry:} header (10-tuple, wildcards) + action + statistics.
\item \textbf{Actions:} forward to port(s); encapsulate to controller (first packet of a flow); drop; plus (for OpenFlow-enabled switches) forward through the normal pipeline. Header rewrite and priority classes are extra (Type 1).
\item \textbf{Isolation:} the fourth action or separate VLANs. Property 1 by the switch, Property 2 by the controller's permissions.
\item \textbf{Controllers:} static (VLAN-like) or dynamic; one switch may have several controllers.
\item \textbf{Flows vs.\ packets:} flow-level decisions are cheap; per-packet processing goes to the controller (slow) or an external NetFPGA (line rate).
\end{itemize}
